\subsection{An explicit finite detector in a one-point Birman kernel}

\begin{result}{A.4.33}{Proved}{Proposition}
\label{prop:birman-cyclic-fibre-detector}
Let $G_p$ be a specified subgroup of the mapping class group of
a punctured sphere which fixes the marked point $p$, and let
\[
 1\longrightarrow F\longrightarrow G_p\longrightarrow G
       \longrightarrow1
\]
be its one-point Birman sequence.  Suppose that the specified
involution $a\in G_p$ is realized by an orientation-preserving
half-turn fixing $p$.  The point-pushing group is the based
free group $F=\pi_1(S,p)$, of rank $r\geq2$, and the actual
conjugation action gives $G_p\to\operatorname{Aut}(F)$,
with $F$ acting by inner automorphisms.
For every specified finite quotient $q_0:G_p\to Q_0$ the
following construction gives a finite quotient $q_+:G_p\to Q_+$
refining $q_0$, with
\[
 q_0\bigl(C_{Q_+}(q_+(a))\cap q_+(F)\bigr)=1.          \tag{BCF.1}
\]
Here the displayed $q_0$ means the retained projection
$Q_+\to Q_0$; thus (BCF.1) is a statement about a literal
finite centralizer fibre.

Put $s=|q_0(F)|$ and define
\[
 J=\kappa_s(F),\qquad J_1=J'J^2,\qquad
 J_2=J_1'J_1^3,\qquad W=F/J_2.                         \tag{BCF.2}
\]
The quotient $q_+$ is the actual diagonal image of $G_p$ under
\[
             q_0\times\bigl(G_p\longrightarrow
                                      \operatorname{Aut}(W)\bigr).
                                                               \tag{BCF.3}
\]
All three subgroups in (BCF.2) are characteristic in $F$.
If $d=[F:J]$, $u=1+d(r-1)$, and $v=d2^u$, then
\[
 d\leq(s!)^{(s!)^r},\qquad
 |W|=v3^{1+v(r-1)},\qquad |Q_+|\leq|Q_0|\,|W|^r.     \tag{BCF.4}
\]
The same finite action separates the two possible order-two
vertex types when the second fixed point of the half-turn is
also unmarked in $S$.
\end{result}

\begin{proof}
There are exactly two fixed points of the specified half-turn
on the sphere.  Suppose $S$ has $n$ deleted points, of which
$f\in\{0,1\}$ are the other fixed point.  The actual orbifold
quotient of $S$ by the half-turn gives the extension
\[
 E=F\rtimes\langle a\rangle
       \cong\mathop{*}_{j=1}^{2-f}C_2
                     *F_{(n+f)/2-1}.                  \tag{BCF.5}
\]
The factor corresponding to the fixed base point $p$ contains
the named element $a$.  The presentation follows by taking
positive loops around the $2-f$ unremoved order-two cone
points and around the $(n+f)/2$ punctures in the quotient
sphere, then eliminating one puncture loop using the product
relation.  The even-cover map sends each order-two factor
to the nonidentity element of $C_2$ and has kernel precisely
$F$.  In particular $C_E(a)=\langle a\rangle$.

The regular representation of $q_0(F)$ is one of the maps
defining $\kappa_s(F)$, so $J\subseteq\ker(q_0|_F)$.
The finite quotient $E/J$ is injective on every order-two
vertex factor, since their nonidentity elements already
map nontrivially to $E/F=C_2$.  Apply the explicit graph
homology argument A.4.7 with the prime two to
$E\to E/J$.  It gives
\[
 C_{E/J_1}(a)\longrightarrow E/J
                     \quad\hbox{with image }\langle a\rangle.
                                                               \tag{BCF.6}
\]
Consequently an element of $F/J_1$ centralizing $a$ has
trivial image in $F/J$.

The extra third-power homology cover in (BCF.2) is needed
because (BCF.3) retains the inner action of $F$, rather than
silently identifying that action with $F/J_1$ itself.
Write $P=F/J_1$.  Its regular graph cover of the rank-$r$
rose has first integral homology of rank $1+|P|(r-1)$.
A nonidentity deck element has integral trace one there,
whereas the identity has trace $1+|P|(r-1)>1$.
The deck action on this homology is therefore faithful.
It remains faithful modulo three: a finite-order integral
matrix congruent to the identity modulo three is the
identity, by the binomial valuation argument (DR.6).
Thus the action of $P$ on $J_1/J_2$ is faithful, and
\[
                    Z(W)\longrightarrow P
                          \quad\hbox{has image }1.     \tag{BCF.7}
\]

Take $f\in F$ whose image under (BCF.3) centralizes the
image of $a$.  On $W$ their commutator is the inner
automorphism defined by $[f,a]\in F$.  Its triviality means
$[f,a]\in Z(W)$.  Formula (BCF.7) places this commutator
in $J_1$; hence $fJ_1$ centralizes $a$ in $E/J_1$.
Formula (BCF.6) puts $f$ in $J$, and therefore $q_0(f)=1$.
This proves (BCF.1).

Here is also the promised vertex-type assertion, including
its precise comparison map.  Let $a_1,a_2$ be the named
nonidentity elements of the two distinct $C_2$ factors
in (BCF.5).  They are not conjugate by the image of $F$
in the finite inner-action representation on $W$.
A claimed conjugacy would, by (BCF.7), give such a
conjugacy in $E/J_1$.  Let $Y=J\backslash T$ be the
finite graph obtained from the Bass--Serre tree of (BCF.5).
A path representing this conjugacy would join a fixed
vertex of type one to a fixed vertex of type two.  Adding
a homology cycle to make the lift commute with the
involution would make that edge-chain invariant.  Every
invariant edge-chain is a sum of two-edge orbits, whose
boundary has coefficient zero at every fixed vertex over
$\mathbb F_2$.  The boundary of the proposed path has
coefficient one at its two distinct endpoint vertices,
a contradiction.  The vertices remain distinct in $Y$
because the quotient by $J$ preserves their vertex types.
This is the same explicit path calculation as A.4.7,
now applied to the conjugacy of two different factors.

The free ranks in (BCF.4) are the Nielsen--Schreier ranks
of $J$ and $J_1$.  They give the displayed orders directly.
An automorphism of the finite group $W$, generated by
$r$ named images, is determined by those $r$ images;
there are at most $|W|^r$ possibilities.  This proves
the final bound and completes the finite construction.
\end{proof}

\paragraph{The six literal free-product inputs in the ten-point marking.}
For the configuration (A.T10), forgetting the ordered fixed
pair $A_0$ one point at a time and using the involution $a$
gives, in order,
\[
 C_2*F_4,\qquad C_2*C_2*F_3.
\]
After that, forgetting $B_0$ with the involution $b$ gives
$C_2*F_3$ and $C_2*C_2*F_2$; finally forgetting $C_0$
with $c=ab$ gives $C_2*F_2$ and $C_2*C_2*F_1$.
The point-pushing free ranks are respectively $8,7,6,5,4,3$.
Each displayed group is the actual orbifold fundamental
group in (BCF.5), with the named involution at the first
fixed-point vertex.  A finite quotient of a point stabilizer
extends to one of the full mixed group by the literal
induced homomorphism into a wreath product over its finite
set of point cosets; an index-$t$ stabilizer and finite
image $Q$ give an image of order at most $t!|Q|^t$.
This records actual finite ambient targets for all six
local detectors.

The proposition identifies the point-pushing part of the
centralizer fibre.  Passing from it to the full simultaneous
$C_2^2$ centralizer still requires the compatible map between
the ordinary centralizer and its forgotten target, including
which fixed-point vertex is lifted.  It is not asserted
here that six unrelated cyclic tests alone establish the
required simultaneous finite comparison.
